\section{Integral discreteness and height gaps}\label{sec:lattices}

\subsection{Real positivity forced by finiteness}\label{sec:northcott-real}

\begin{lemma}\label[lemma]{lem:northcott-real}
Let \(\Lambda\) be a subgroup spanning a rational vector space
\(V\), and let \(q\) be a real-valued quadratic form on \(V\).
If \(q\) has Northcott on \(\Lambda\), its extension to
\(V\otimes_\Q\R\) is positive definite.
\end{lemma}

\begin{proof}
It suffices to consider a finite-dimensional rational subspace
and choose a rational basis \(v_1,\ldots,v_s\) in \(\Lambda\).
Let \(Q\) be the Gram matrix. A negative real direction would,
by rational approximation and clearing denominators, give a
nonzero \(v\in\sum_i\Z v_i\) with \(q(v)<0\). Its integer
multiples contradict Northcott. Thus \(Q\) is positive semidefinite.

If \(0\ne u\in\ker Q\), round the coordinates of \(tu\) to
integers, obtaining \(z_t\in\Z^s\) with uniformly bounded
\(z_t-tu\). Then
\[
 q\Bigl(\sum_i(z_t)_iv_i\Bigr)
       =(z_t-tu)^tQ(z_t-tu)
\]
is uniformly bounded. The rounded vectors are unbounded as
\(t\to\infty\), hence take infinitely many distinct values.
This again contradicts Northcott. Therefore every such finite
Gram matrix is positive definite, which proves the assertion.
\end{proof}

In particular, residual integral Northcott already contains the
real-positivity assertion for \(q_G\). The converse requires an
arithmetic condition on the integral image.

\subsection{Finite-rank lattices}\label{sec:finite-rank-lattices}

\begin{theorem}\label[theorem]{thm:lattice-equivalence}
Let \(\Lambda\) be torsion-free of finite rank,
\(V=\Lambda\otimes\Q\), and suppose that \(q\) is real positive
definite on \(V\otimes\R\). The following are equivalent:
\begin{enumerate}[label=(\roman*)]
\item \(q\) has Northcott on \(\Lambda\).
\item There is \(\varepsilon>0\) such that
\(q(\lambda)\ge\varepsilon\) for every nonzero \(\lambda\in\Lambda\).
\item \(\Lambda\) is finitely generated.
\item For a rational basis \(v_1,\ldots,v_s\) of \(V\), there is
an integer \(N\ge1\) such that
\(\Lambda\subset N^{-1}\sum_i\Z v_i\).
\item The diagonal image of \(\Lambda\) in \(V(\A_f)\) is
relatively compact.
\end{enumerate}
The zero group satisfies these statements with the gap condition
interpreted vacuously.
\end{theorem}

\begin{proof}
Finiteness in a unit height ball gives a positive minimum, or the
bound \(1\) if that ball contains no nonzero point. This proves
(i) implies (ii). The gap makes \(\Lambda\) uniformly discrete
for the Euclidean norm \(\sqrt q\). Choose a rational basis in
\(\Lambda\) and reduce elements of \(\Lambda\) modulo its
\(\Z\)-span into a bounded fundamental parallelepiped. Uniform
discreteness makes the set of representatives finite. Hence this
span has finite index in \(\Lambda\), proving (iii) and (iv).
Conversely, a bounded subset of a Euclidean lattice is finite,
so (iv) implies (i). Finite generation gives (iv) for any rational
basis by clearing the finitely many coefficients.

Every compact subset of \(V(\A_f)\) is contained in
\(N^{-1}\Lambda_0\otimes\widehat\Z\) for some integer \(N\)
and a fixed reference lattice \(\Lambda_0\). This follows by a
finite subcover from the open compact sets of that form. Conversely
each such set is compact. Intersecting with \(V(\Q)\) proves the
equivalence with (v).
\end{proof}

If \(\Lambda\subset N^{-1}\Lambda_0\), and the least eigenvalue
of the Gram matrix in a basis of \(\Lambda_0\) is \(\mu>0\),
one obtains the explicit estimate
\begin{equation}\label{eq:lattice-count}
 \#\{\lambda\in\Lambda:q(\lambda)\le B\}
 \le\left(1+2\left\lfloor N\sqrt{B/\mu}\right\rfloor\right)^s
 \qquad(B\ge0).
\end{equation}

\subsection{An all-prime saturation criterion}\label{sec:all-prime}

\begin{theorem}\label[theorem]{thm:all-prime}
Let \(V\) be a finite-dimensional rational vector space,
\(\Lambda_0\subset V\) a full lattice, and \(\Lambda\subset V\)
a subgroup. For every prime \(p\), let
\[
 D_p:V\otimes\Q_p\hookrightarrow W_p
\]
be an injective linear map to a finite-dimensional space, and let
\(M_p\subset W_p\) be a full \(\Z_p\)-lattice. Assume
\[
 D_p(\Lambda)\subseteq M_p\quad\text{for every }p,
 \qquad D_p^{-1}(M_p)=\Lambda_0\otimes\Z_p
               \quad\text{for almost every }p.
\]
Then \(\Lambda\subset N^{-1}\Lambda_0\) for one integer \(N\).
Every real positive definite quadratic form has Northcott on
\(\Lambda\).
\end{theorem}

\begin{proof}
Injectivity makes \(L_p=D_p^{-1}(M_p)\) a \(\Z_p\)-lattice:
it is open and compact in \(V\otimes\Q_p\). For each exceptional
prime choose \(b_p\ge0\) with
\(L_p\subset p^{-b_p}\Lambda_0\otimes\Z_p\), and put
\(N=\prod p^{b_p}\). The coordinates of \(N\lambda\) lie in
\(\Z_p\) for every \(p\). Since
\(\Q\cap\prod_p\Z_p=\Z\), they are integers. Apply
\Cref{thm:lattice-equivalence}.
\end{proof}

For the residual regulator, a common multiple of the projector
denominators can first be absorbed into \(D_p\). Once scalar
injectivity is available, \Cref{prop:local-lattice} supplies local
integral modules. The almost-everywhere saturation condition in
\Cref{thm:all-prime} is the additional comparison with one common
rational lattice.

\begin{example}\label[example]{ex:denominators}
Let
\[
 \Lambda=\sum_{p\ {
m prime}}\Z\,\frac1p\subset\Q,
 \qquad q(x)=x^2.
\]
Its elements have squarefree denominators. For every prime \(p\),
\(\Lambda\otimes\Z_p=p^{-1}\Z_p\), which is finite free of
rank one. Nevertheless \(q(1/p)=p^{-2}\to0\), so Northcott
fails. The integral maps \(D_p(x)=px\) have inverse lattices
larger than \(\Z_p\) at infinitely many primes. Thus local
finite generation at every prime does not provide the common
denominator required by \Cref{thm:lattice-equivalence}.
\end{example}

Already \(\Z[1/p]\), with \(q(x)=x^2\), defeats any finite
collection of integral detectors at primes different from \(p\).
Another distinction is real definiteness: the form
\((a-\sqrt2b)^2\) on \(\Z^2\) is positive on nonzero rational
vectors, but rounding along its real kernel produces infinitely
many bounded-height lattice vectors.

\subsection{A quantitative descent criterion}\label{sec:weak-mw}

\begin{theorem}\label[theorem]{thm:weak-mw-gap}
Let \(\Lambda\) be torsion-free and let \(\sqrt q\) be the norm
of a real positive quadratic form on \(\Lambda\otimes\R\).
Suppose, for some integer \(m\ge2\),
\[
 s=|\Lambda/m\Lambda|<\infty,\qquad
 q(\lambda)\ge\varepsilon>0\quad(0\ne\lambda\in\Lambda).
\]
Then \(\Lambda\) is finitely generated and \(q\) has Northcott.
More precisely, if \(B\ge0\), \(j\in\Z_{\ge0}\), and
\(m^{2j}\varepsilon>4B\), then
\begin{equation}\label{eq:kummer-count}
 \#\{\lambda\in\Lambda:q(\lambda)\le B\}\le s^j.
\end{equation}
\end{theorem}

\begin{proof}
If \(x,y\) are in the height-\(B\) ball and
\(x-y=m^jz\), then
\(\|z\|\le2\sqrt B/m^j<\sqrt\varepsilon\), so \(z=0\).
Reduction modulo \(m^j\Lambda\) is injective on that ball.
Torsion-freeness identifies the successive quotients
\(m^i\Lambda/m^{i+1}\Lambda\) with \(\Lambda/m\Lambda\),
proving \MathEqref{eq:kummer-count}.

Choose a finite set \(T\) of representatives modulo
\(m\Lambda\), and let \(C=\max_{t\in T}\|t\|\).
Writing \(x=my+t\) gives
\(\|y\|\le(\|x\|+C)/m\). Repetition brings \(y\) into
any fixed ball of radius greater than \(C/(m-1)\). This ball
is finite by the first part; together with \(T\), it generates
\(\Lambda\). This is the height-descent argument underlying
weak Mordell--Weil; compare \cite[Section~10.3]{ConradFeng}.
\end{proof}

Northcott by itself need not give finite rational rank. For example,
\(\Lambda=\bigoplus_{i\ge1}\Z e_i\) with
\(q(\sum a_ie_i)=\sum2^ia_i^2\) has finite height balls and
infinite rank. Rational detection, integral discreteness and real
positivity therefore retain distinct roles in the general program.

The examples in this section are examples of abstract groups and
quadratic forms. Their purpose is to identify the exact arithmetic
hypotheses in a proof of geometric Northcott.
